\documentclass[12pt]{article}
\usepackage{amsmath, amssymb}
\usepackage{geometry}
\geometry{margin=1in}
\usepackage{setspace}
\setstretch{1.2}

\title{CSE400 – Fundamentals of Probability in Computing}
\author{}
\date{}

\begin{document}

\maketitle

\section*{Lecture 9: Uniform and Exponential Random Variables}

\hrule
\vspace{0.5cm}

\section*{1. Types of Continuous Random Variables (Overview)}

The lecture introduces the following continuous random variables:

\begin{itemize}
\item Uniform Random Variable
\item Exponential Random Variable
\item Laplace Random Variable (listed but not developed in shown portion)
\item Gamma Random Variable (listed but not developed in shown portion)
\end{itemize}

The focus in this material is on:

\begin{itemize}
\item Uniform Random Variable
\item Exponential Random Variable
\end{itemize}

\section*{2. Uniform Random Variable}

Let \( X \) be a continuous random variable uniformly distributed over the interval \( [a, b] \).

\subsection*{2.1 Probability Density Function (PDF)}

\[
f_X(x) =
\begin{cases}
\dfrac{1}{b-a}, & a \le x < b \\
0, & \text{elsewhere}
\end{cases}
\]

This means:

\begin{itemize}
\item The density is constant between \( a \) and \( b \)
\item Zero outside the interval
\end{itemize}

\subsection*{2.2 Cumulative Distribution Function (CDF)}

\[
F_X(x) =
\begin{cases}
0, & x < a \\
\dfrac{x-a}{b-a}, & a \le x < b \\
1, & x \ge b
\end{cases}
\]

Thus:

\begin{itemize}
\item Before \( a \), probability is zero
\item Between \( a \) and \( b \), probability increases linearly
\item After \( b \), probability equals one
\end{itemize}

\subsection*{2.3 General Probability for Uniform RV on \( [0, 2\pi] \)}

For a uniform random variable \( \Theta \in [0,2\pi] \):

\[
\Pr(a < \Theta < b) = \frac{b-a}{2\pi}
\]

\section*{3. Example 1 — Uniform Random Variable}

\subsection*{Given:}

The phase \( \Theta \) of a sinusoid is uniformly distributed over \( [0,2\pi] \):

\[
f_\Theta(\theta) =
\begin{cases}
\dfrac{1}{2\pi}, & 0 \le \theta < 2\pi \\
0, & \text{otherwise}
\end{cases}
\]

\subsection*{(a) Find \( \Pr(\Theta > \frac{3\pi}{4}) \)}

Using uniform probability:

\[
\Pr(\Theta > \frac{3\pi}{4})
= \frac{2\pi - \frac{3\pi}{4}}{2\pi}
\]

\[
= \frac{8\pi - 3\pi}{8\pi}
= \frac{5\pi}{8\pi}
= \frac{5}{8}
\]

\subsection*{(b) Find \( \Pr(\Theta < \pi \mid \Theta > \frac{3\pi}{4}) \)}

Use conditional probability:

\[
\Pr(A|B) = \frac{\Pr(A \cap B)}{\Pr(B)}
\]

Here:

\begin{itemize}
\item \( B = \Theta > \frac{3\pi}{4} \)
\item \( A = \Theta < \pi \)
\end{itemize}

First compute:

\[
\Pr\left(\frac{3\pi}{4} < \Theta < \pi\right)
= \frac{\pi - \frac{3\pi}{4}}{2\pi}
= \frac{\pi/4}{2\pi}
= \frac{1}{8}
\]

From part (a):

\[
\Pr(B) = \frac{5}{8}
\]

Therefore:

\[
\Pr(\Theta < \pi \mid \Theta > \frac{3\pi}{4})
= \frac{1/8}{5/8}
= \frac{1}{5}
\]

\subsection*{(c) Find \( \Pr(\cos\Theta < \frac{1}{2}) \)}

Solve:

\[
\cos\Theta = \frac{1}{2}
\Rightarrow \Theta = \frac{\pi}{3}, \frac{5\pi}{3}
\]

Thus:

\[
\cos\Theta < \frac{1}{2}
\quad \text{for} \quad \frac{\pi}{3} < \Theta < \frac{5\pi}{3}
\]

Probability:

\[
\Pr = \frac{\frac{5\pi}{3} - \frac{\pi}{3}}{2\pi}
= \frac{4\pi/3}{2\pi}
= \frac{2}{3}
\]

\section*{4. Applications of Uniform Random Variable}

Examples stated:

\begin{itemize}
\item Phase of a sinusoidal signal uniformly between \( 0 \) and \( 2\pi \)
\item Random number between 0 and 1 in simulations
\item Arrival time within a known time window assuming no preference
\end{itemize}

\section*{5. Exponential Random Variable}

Let \( b > 0 \).

\subsection*{5.1 Probability Density Function (PDF)}

\[
f_X(x) = \frac{1}{b} \exp\left(-\frac{x}{b}\right) u(x)
\]

where \( u(x) \) is the unit step function ensuring:

\[
x \ge 0
\]

\subsection*{5.2 Cumulative Distribution Function (CDF)}

\[
F_X(x) = \left[1 - \exp\left(-\frac{x}{b}\right)\right] u(x)
\]

Thus:

\begin{itemize}
\item Probability starts at zero
\item Increases exponentially toward 1 as \( x \to \infty \)
\end{itemize}

\section*{6. Example 2 — Exponential Random Variable}

\subsection*{Given:}

\[
f_X(x) = e^{-x} u(x)
\]

This corresponds to exponential RV with \( b = 1 \).

\subsection*{(a) Find \( \Pr(3X < 5) \)}

First rewrite inequality:

\[
3X < 5 \Rightarrow X < \frac{5}{3}
\]

Now use CDF:

\[
\Pr(X < y) = 1 - e^{-y}
\]

So:

\[
\Pr(3X < 5) = 1 - e^{-5/3}
\]

\subsection*{(b) Generalize: Find \( \Pr(3X < y) \)}

Rewrite:

\[
3X < y \Rightarrow X < \frac{y}{3}
\]

Thus:

\[
\Pr(3X < y) = 1 - e^{-y/3}
\]

\end{document}
